\documentclass[11pt]{article}

\usepackage[T1]{fontenc}
\usepackage[utf8]{inputenc}
\usepackage{lmodern}
\usepackage[margin=1in]{geometry}
\usepackage{microtype}
\usepackage{amsmath}
\usepackage{booktabs}
\usepackage{longtable}
\usepackage{tabularx}
\usepackage{array}
\usepackage{xcolor}
\usepackage{enumitem}
\usepackage{listings}
\usepackage{fancyhdr}
\usepackage[hidelinks]{hyperref}

% Keep the generated PDF stable across identical builds with pdfTeX.
\pdfinfoomitdate=1
\pdftrailerid{}
\pdfsuppressptexinfo=15

\definecolor{pipelineblue}{HTML}{315E7D}
\definecolor{lightblue}{HTML}{EDF4F8}
\definecolor{lightgold}{HTML}{FFF7E6}
\definecolor{codegray}{HTML}{F4F5F6}

\hypersetup{
  pdftitle={Drosophila cCRE pipeline: detailed implementation},
  pdfauthor={drosophila-cCREs project},
  pdfsubject={ATAC-seq, H3K27ac ChIP-seq, master DHS, and regulatory catalog workflow}
}

\pagestyle{fancy}
\fancyhf{}
\fancyhead[L]{Drosophila cCRE pipeline}
\fancyhead[R]{Detailed implementation}
\fancyfoot[C]{\thepage}
\setlength{\headheight}{14pt}

\setlist[itemize]{leftmargin=1.6em,itemsep=0.25em,topsep=0.35em}
\setlist[enumerate]{leftmargin=1.8em,itemsep=0.3em,topsep=0.35em}
\setlength{\parindent}{0pt}
\setlength{\parskip}{0.55em}
\setlength{\emergencystretch}{3em}
\renewcommand{\arraystretch}{1.18}

\lstdefinestyle{pipelinecode}{
  basicstyle=\ttfamily\small,
  backgroundcolor=\color{codegray},
  frame=single,
  rulecolor=\color{gray!35},
  breaklines=true,
  columns=fullflexible,
  keepspaces=true,
  showstringspaces=false,
  xleftmargin=0.3em,
  xrightmargin=0.3em
}
\lstset{style=pipelinecode}

\newcommand{\stage}[1]{\texttt{#1}}
\newcommand{\filep}[1]{\path{#1}}
\newcommand{\software}[1]{\textsf{#1}}

\title{\textbf{Drosophila cCRE Pipeline}\\[0.3em]
\large Detailed implementation and scientific reasoning}
\author{drosophila-cCREs project}
\date{Repository documentation revision: 2026-07-29}

\begin{document}

\maketitle

\begin{abstract}
This document describes the implemented workflow that converts public
ATAC-seq and ChIP-seq accessions into quality-reviewed biological-library
BAMs, replicate-supported context peaks, a variable-width summit-aware master
DNase-hypersensitive-site (DHS) registry, and a context-resolved regulatory
element catalog.  The downstream catalog combines ATAC accessibility and
H3K27ac signal using assay-specific background TMM normalization, guarded
two-component Gaussian mixture models, portable browser tracks, IGV sessions,
and an integrated checksummed report.  The emphasis is both operational and
scientific: each section states what the current rules do, why that choice was
made, and which limitations remain.  This is implementation documentation for
the live \filep{workflow/Snakefile} and its included rule files; it is not a
substitute for inspecting the resolved configuration and provenance of a
particular run.
\end{abstract}

\tableofcontents
\newpage

\section{Purpose and scope}

The pipeline has two related scientific goals:

\begin{enumerate}
  \item construct a stable atlas coordinate system---the master DHS registry---
  from reproducible ATAC accessibility evidence across the contexts named in
  the input table; and
  \item annotate every master DHS in every context with accessibility,
  H3K27ac, an activity score, TSS proximity, and a guarded active/inactive
  interpretation.
\end{enumerate}

The public interface is accession-first.  Users provide public accessions and
biological metadata, not raw FASTQ paths.  The canonical launcher is
\filep{src/run_pipeline.py}; it validates inputs, resolves defaults and
accession layouts, writes a deterministic YAML configuration, and invokes the
canonical \filep{workflow/Snakefile}.  Snakemake then owns dependency tracking,
resource scheduling, rule environments, and restart behavior.

The current workflow includes ATAC-seq processing, TF and histone ChIP-seq
processing, H3K27ac-based context activity annotation, catalog visualization,
and integrated reporting.  It does not perform ABC candidate resizing or ABC
scoring, and it does not integrate Micro-C contacts.  It also does not import
an external regulatory-element catalog or apply quantile normalization.

A master registry should be derived from the project's own reviewed
accessibility evidence.  Keeping contact-based enhancer--gene assignment and
ABC scoring outside this workflow prevents a coordinate-building pipeline from
silently making a different biological claim about target genes.

\section{The data model}

\subsection{Run, library, context, and assay}

Four identifiers have distinct meanings:

\begin{description}[style=nextline,leftmargin=2.2em]
  \item[Accession or technical run]
  One public run, or one run obtained by expanding a public experiment
  accession.  Technical runs are processed independently through trimming and
  alignment.

  \item[Biological library]
  The unit identified by \texttt{library\_id}.  Multiple technical runs may
  share one library ID.  Their aligned BAMs are merged before duplicate
  marking so duplicates are defined across the complete biological library.

  \item[Context]
  The tissue, stage, cell type, or other biological group in the sample-sheet
  \texttt{context} column.  Distinct ATAC libraries in one context are
  biological replicates.  The same context label also links accepted ATAC and
  H3K27ac libraries during activity quantification.

  \item[Assay]
  One of \texttt{atac}, \texttt{h3k27ac}, \texttt{chip\_tf}, or
  \texttt{chip\_histone}.  The \texttt{h3k27ac} label resolves to the histone
  ChIP processing branch but is retained as H3K27ac in the downstream atlas.
\end{description}

Collapsing these levels would create pseudoreplication or biased weights.  A
technical lane is not a biological replicate; a library should receive one
QC decision; and a context estimate should give equal weight to accepted
biological libraries rather than to the number of sequencing lanes.

\subsection{Canonical sample sheet}

The CSV or TSV input conforms to
\filep{schemas/sample-sheet.schema.yaml}.  Four columns are required:

\begin{center}
\begin{tabularx}{0.94\textwidth}{>{\ttfamily}l X}
\toprule
\normalfont Column & Meaning \\
\midrule
accession & Public run or experiment accession; experiments expand to their runs. \\
library\_id & Biological library identifier; shared by technical runs only. \\
assay & Assay branch and its default peak behavior. \\
context & Replicate-pooling and downstream atlas group. \\
\bottomrule
\end{tabularx}
\end{center}

Optional fields specify treatment/control roles, an explicit
\texttt{control\_library}, MACS3 peak-calling parameters, adapters,
alignment parameters, filtering parameters, and notes.  ChIP input controls
are explicit relationships.  A blank control is a valid IP-only experiment
and causes MACS3 to run without \texttt{-c}; it does not trigger an implicit
control search.

Genome is a project-level launcher argument (\texttt{dm6} or \texttt{hg38}),
and sequencing layout is resolved from accession metadata.  Repeated
per-row genome or raw FASTQ fields are deliberately absent from the public
contract.

\section{Staged execution and the manual checkpoint}

\subsection{Why there are separate invocations}

The workflow is intentionally not one uninterrupted accession-to-catalog job.
It is organized around immutable, checksummed artifacts and an explicit human
review boundary:

\begin{center}
\fcolorbox{pipelineblue}{lightblue}{%
\begin{minipage}{0.92\textwidth}
\textbf{Invocation 1: accessions $\rightarrow$ QC}\\
download $\rightarrow$ trim $\rightarrow$ align $\rightarrow$ peaks and QC
$\rightarrow$ pending final-BAM manifest and review table

\vspace{0.35em}
\centering $\Downarrow$ \quad \textbf{manual pass/fail review} \quad $\Downarrow$

\raggedright\vspace{0.35em}
\textbf{Invocation 2: reviewed final BAMs $\rightarrow$ master}\\
validate $\rightarrow$ accepted ATAC replicate/context peaks $\rightarrow$
master DHS bundle

\vspace{0.35em}
\textbf{Invocation 3: accepted BAMs + frozen master $\rightarrow$ catalog/report}\\
validate $\rightarrow$ count $\rightarrow$ normalize $\rightarrow$ annotate
$\rightarrow$ tracks, IGV sessions, and report
\end{minipage}}
\end{center}

There is no interactive pause inside Snakemake.  The user copies the generated
review table outside the result namespace, records one pass/fail decision for
every library, and applies those decisions with
\filep{src/review_final_bam_manifest.py}.  A master run rejects any pending
ATAC decision.  Quantification similarly requires accepted ATAC and H3K27ac
artifacts; a rejected row must carry a reason.

An automatic pipeline can calculate metrics but cannot decide whether a
library is biologically fit for a specific atlas.  Making the review boundary
an explicit file transformation preserves human judgment without hiding it in
an interactive process or modifying generated manifests in place.

\subsection{Input and stopping stages}

\begin{longtable}{>{\ttfamily}p{0.18\textwidth} p{0.34\textwidth} p{0.35\textwidth}}
\toprule
\normalfont Start & Required starting artifact & Allowed endpoints \\
\midrule
\endhead
accessions & Sample sheet plus resolved, checksummed download manifest &
\stage{trimming}, \stage{alignment}, or \stage{qc} \\
trimming & Checksummed trimming checkpoint &
\stage{trimming}, \stage{alignment}, or \stage{qc} in the source namespace \\
alignment & Checksummed alignment checkpoint or complete final-BAM manifest &
\stage{alignment} or \stage{qc} \\
qc & QC checkpoint plus a reviewed ATAC final-BAM manifest for master construction &
\stage{qc} or \stage{master} \\
master & Checksummed master bundle and accepted ATAC/H3K27ac BAM manifests when continuing &
\stage{master}, \stage{quantification}, \stage{catalog}, or \stage{report} \\
quantification & Checksummed quantification checkpoint &
\stage{quantification}, \stage{catalog}, or \stage{report} \\
catalog & Checksummed catalog checkpoint & \stage{catalog} or \stage{report} \\
report & Checksummed report checkpoint & \stage{report} validation \\
\bottomrule
\end{longtable}

Each endpoint exports
\filep{provenance/manifests/<stage>.checkpoint.json}.  It records the stage,
scientific semantic digest, complete peak/refinement parameters where relevant,
and SHA-256 hashes of the restart artifacts.  Its stage-specific resolved
configuration snapshot is retained below \filep{provenance/configs/}, so
advancing the same run does not invalidate an earlier boundary.  Advancing
\texttt{--until-stage} inside an allowed start mode changes only requested
targets; it is excluded from the scientific semantic digest.

Strict reuse modes have no upstream fallback.  For example, a quantification
run cannot silently redownload or realign a missing library.  It first verifies
the declared files, SHA-256 hashes, genome, BAM/index coherence, coordinate
compatibility, and filtering contract.

\section{Reference preparation}

For \texttt{dm6} and \texttt{hg38}, reference construction belongs to the
workflow DAG.  The resolved configuration records canonical FASTA, NCBI
RefSeq GTF, and Blacklist v2 sources together with MD5 or SHA-256 checksums.
The workflow:

\begin{enumerate}
  \item downloads archives with resumable \software{aria2} transfers and
  integrity checking;
  \item decompresses the FASTA and blacklist through temporary files followed
  by atomic replacement;
  \item builds the \software{samtools faidx} and \software{Bowtie2} indexes;
  \item derives chromosome sizes in FASTA order;
  \item extracts one-base TSS records from the annotation; and
  \item writes the configured autosome list used for background normalization.
\end{enumerate}

The blacklist is used during final BAM filtering.  The TSS BED supports ATAC
TSS enrichment and regulatory-class annotation.  The
autosome list excludes sex chromosomes and mitochondrial DNA from the TMM
background matrix.

Reference identity is part of the computation, not a machine-local
prerequisite.  Checksummed preparation in the DAG makes alignment, QC,
normalization, and coordinates traceable to the same assembly and avoids
silent reuse of a similarly named but different reference.

\section{Acquisition, read QC, and trimming}

\subsection{Checksummed acquisition}

Run and experiment accessions are normalized and experiments are expanded to
their constituent runs.  Acquisition produces immutable compressed FASTQs and
a manifest with resolved paths, layout, and checksums.  Resume state is
preserved, and SRA-derived files are promoted only after conversion and
compression complete.  The \texttt{--skip-download} path reuses an existing
manifest after validation.

\subsection{Raw and trimmed read assessment}

Every technical run and mate receives raw \software{FastQC}.  Reads are then
trimmed with \software{Cutadapt}; defaults are Nextera adapters for ATAC and
TruSeq adapters for ChIP, quality cutoff 20, minimum retained length 20 bp,
error rate 0.1, and minimum adapter overlap 3 bp.  Paired-end trimming is
synchronized.  FastQC is repeated on the trimmed output, and Cutadapt writes a
JSON record used by MultiQC.

Running QC before and after trimming separates properties of the deposited
reads from effects introduced by processing.  Keeping technical runs separate
at this stage reveals lane-specific failures before they are hidden by a
library merge.

\section{Alignment and biological-library BAM construction}

\subsection{Lane alignment}

Each trimmed technical run is aligned independently with \software{Bowtie2}.
The default preset is \texttt{--very-sensitive}.  For paired-end data, the
workflow also uses \texttt{--no-mixed}, \texttt{--no-discordant}, and a default
maximum fragment length of 2,000 bp.  Read-group IDs identify the technical
run while the sample field identifies the biological library.

The streamed alignment is converted to BAM, name-sorted, processed with
\software{samtools fixmate -m}, and coordinate-sorted.  Temporary intermediates
are checked with \software{samtools quickcheck}.

\subsection{Merge, duplicate marking, and filtering}

All lane BAMs sharing a biological library ID are merged.  Duplicate marking
is then performed once over the merged library.  The final filter applies:

\begin{itemize}
  \item minimum MAPQ 30 by default;
  \item exclusion of unmapped, secondary, QC-failed, and supplementary
  alignments;
  \item proper-pair and mapped-mate requirements for paired-end libraries;
  \item duplicate exclusion when \texttt{remove\_duplicates=true};
  \item mitochondrial-contig removal by default; and
  \item removal of alignments overlapping the reference blacklist.
\end{itemize}

The result is a coordinate-sorted final BAM and BAI.  Flagstat, stats, and
idxstats summaries are calculated for every biological library.  At the
\stage{alignment} boundary the workflow exports a checksummed final-BAM
manifest; newly generated rows have \texttt{qc\_status=pending\_review}.

Duplicate marking before technical-run merging would miss molecules duplicated
across lanes.  Conversely, merging biological replicates would erase the
replicate structure required for reproducibility checks.  The chosen order
therefore matches the experimental unit.

\section{Peak calling}

\subsection{Default paired-end ATAC branch: Tn5/MACS3-qpois}

The default ATAC implementation is explicitly insertion-based:

\begin{enumerate}
  \item Retain proper pairs with nonzero template length and
  \(0 < |\mathrm{TLEN}| < 150\) bp by default.  Unmapped, mate-unmapped,
  secondary, QC-failed, duplicate, and supplementary records are excluded.
  \item Apply the Tn5 offset with \software{alignmentSieve --ATACshift}.
  \item Convert both shifted mates to separate one-base insertion records.  A
  consistency check requires the shifted alignment count to equal the number
  of insertion records.
  \item Run \software{MACS3 callpeak} on the BED input with
  \texttt{-f BED --nomodel --shift -75 --extsize 150 --keep-dup all -B
  -q 0.10}.  The pileup and local lambda are deliberately unscaled; this
  branch must not use \texttt{--SPMR}.
  \item Clip the treatment pileup and local lambda to chromosome bounds and
  calculate the qpois signal with \software{macs3 bdgcmp -m qpois}.
  \item Refine candidate peaks over qpois exponents 2 through 325.  Connected
  components shorter than 50 bp are rejected; components at most 400 bp are
  retained; wider components are followed to higher exponents so that they can
  split into narrower subcomponents.  The default merge gap is one base.
\end{enumerate}

This sequence runs for each accepted biological replicate and again on pooled
insertions within each context.  The pooled call provides context boundaries
and signal; the replicate calls provide reproducibility evidence.
The QC endpoint freezes each refined replicate BED together with the complete
MACS3 and qpois-refinement parameter mappings.  A QC-to-master invocation
validates and reuses those lenient replicate calls, while still constructing
the pooled context calls.  It refuses reuse if the requested q-value, Tn5
geometry, exponent range, component lengths, or merge gap differs.

Two-ended Tn5 insertions represent cut sites more directly than treating an
ATAC fragment as a ChIP enrichment interval.  Computing qpois from unscaled
pileup and lambda preserves MACS3's intended count-scale comparison.
Progressive refinement constrains the geometry of broad candidates, but it is
not an additional independent FDR-controlled test; candidate q-values remain
the statistical starting point.

\subsection{ChIP branches}

ChIP treatment libraries use MACS3 with their explicit control BAM when one is
declared and without \texttt{-c} for IP-only data.  TF ChIP defaults to narrow
peaks; histone ChIP, including H3K27ac, defaults to broad peaks.  The default
q-value is 0.01 and the default broad cutoff is 0.1.  MACS3 format resolves
from layout unless explicitly constrained.

Both narrow and broad ChIP calls request \texttt{-B --SPMR}, producing
treatment-pileup and control-lambda bedGraphs in addition to the canonical
narrowPeak or broadPeak file.  Per-library ChIP CPM BigWigs are not generated
because no downstream endpoint consumes them.  The SPMR behavior is
intentionally different from the ATAC qpois branch.

Narrow TF occupancy and broad histone enrichment have different expected
geometries.  Explicit control relationships prevent accidental cross-context
matching, while allowing genuinely IP-only public experiments to remain
representable and auditable.

\section{Quality control and manual acceptance}

The QC endpoint combines general and assay-specific evidence:

\begin{itemize}
  \item raw/trimmed FastQC and Cutadapt metrics;
  \item final-BAM flagstat, stats, and idxstats;
  \item peak count and fraction of reads in peaks (FRiP);
  \item ATAC TSS profiles over \(\pm2\) kb, derived from temporary Tn5-shifted
  CPM tracks, and fragment-length histograms;
  \item ChIP treatment/control fingerprints when a control exists;
  \item phantompeakqualtools cross-correlation, fragment-shift estimates,
  NSC, RSC, and quality tags for ChIP treatments; and
  \item a MultiQC report and one consolidated review row per library.
\end{itemize}

For paired-end libraries, FRiP uses one valid same-chromosome fragment per read
pair in the denominator and counts a fragment once if it overlaps any called
peak.  For single-end libraries, the unit is one filtered read.  Thus FRiP is
fragment-aware and is not inflated by counting two mates independently.

ATAC TSS enrichment is calculated from the aggregate 4-kb profile as the
maximum signal in the central \(\pm50\)-bp interval divided by the mean signal
in the two terminal 100-bp flanks.  The review table also reports the median
ATAC fragment length, fraction below 150 bp, and fraction from 180--250 bp.

The user edits only \texttt{qc\_decision},
\texttt{estimated\_fragment\_length\_bp}, and \texttt{notes}.  A failed row
requires a reason.  Single-end H3K27ac also requires a reviewed positive
fragment length; the cross-correlation estimate is pre-populated as a
suggestion rather than accepted silently.

No universal numeric cutoff is safe across all public datasets.  The table
puts depth, mapping, enrichment, fragment structure, and visual evidence next
to the decision so inclusion criteria can be documented at library resolution.

\section{Context peak reproducibility}

Accepted ATAC libraries are grouped by the sample-sheet context.  Their
insertion BEDs are concatenated and a pooled qpois context call is made and
refined.  The pooled peaks are not accepted merely because they are present in
the aggregate signal.

For every pooled peak \(P\), each replicate peak set is reduced to an interval
union and the coverage fraction is calculated:

\[
  f_r(P) = \frac{|P \cap U_r|}{|P|},
\]

where \(U_r\) is the union of replicate \(r\)'s peaks.  A replicate supports
the pooled peak when \(f_r(P) \geq 0.5\) by default.  The peak is retained when
at least two biological libraries support it.  Both thresholds are launcher
parameters, and the minimum cannot exceed the accepted replicate count.  The
support table retains every pooled peak, all per-replicate fractions, the
supporting library list, and the retain/reject decision.

Pooled data improve sensitivity and define a context-level signal, but pooling
alone can promote a feature driven by one library.  Requiring substantial
coverage in multiple independent biological libraries makes reproducibility
an explicit property of every retained context peak.

\section{Summit-aware master DHS registry}

The master registry is built only from replicate-supported context peaks.  It
is neither a simple interval union nor a fixed-width resizing operation.

\subsection{Context summits}

For each source peak, the workflow reads the pooled context BigWig and chooses
the maximum-signal base.  If the maximum occupies more than one plateau, it
selects the midpoint of the plateau closest to the peak midpoint, with a
leftmost deterministic tie break.  If no finite signal is available, it uses
the peak midpoint and records a summit fallback.

\subsection{Cross-context clustering}

Source peaks are processed by chromosome.  A peak can join a cluster only if:

\begin{itemize}
  \item its context is not already represented in that cluster;
  \item the complete span of member summits remains at most 150 bp by default;
  and
  \item it reciprocally contains a summit with at least one current member:
  each of the two peak intervals contains the other's summit.
\end{itemize}

When multiple clusters are eligible, distance to the cluster median summit,
then interval-envelope width, gives a deterministic choice.  The representative
summit is the observed member summit closest to the median; ties prefer a
narrower source peak, then stable context/name order.

Adjacent clusters whose representative summits are less than 50 bp apart are
merged only when their context sets are disjoint and the combined summit span
still satisfies the 150-bp limit.  A shared context is evidence that the same
biological condition resolves two sites, so those sites remain separate.

Each master interval initially spans the envelopes of its members.  If adjacent
master intervals overlap, their boundary is clipped at the midpoint between
their representative summits, preserving both summits and yielding sorted,
non-overlapping, variable-width BED6 intervals.

\subsection{Master outputs}

The bundle contains the master BED, one-base summit BED, full source-peak
membership table, binary context matrix, and JSON metrics.  The membership
table records input and clipped coordinates, source summit and signal,
fallback status, peak method, replicate support, and supporting libraries.  A
manifest records paths and SHA-256 hashes for the complete bundle.

The summit is a sharper cross-context anchor than a peak edge, while preserving
source envelopes avoids pretending that every DHS has the same width.
Reciprocal summit containment reduces accidental joins between nearby broad
peaks.  The one-peak-per-context constraint and shared-context separation rule
protect biologically resolved neighboring sites from being collapsed.

\section{Strict artifact reuse and validation}

Final-BAM and master manifests are immutable contracts between phases.  Reuse
validates file presence, SHA-256, BAM/BAI consistency, genome, coordinate
compatibility, expected filtering metadata, QC status, and master bundle
structure before downstream work.  Rejected libraries remain visible in the
reviewed manifest and exclusion provenance, but they do not enter master or
activity calculations.

The result namespace \filep{results/<project>/<run_id>} is protected by a
timestamp-independent scientific semantic digest.  If an existing namespace
contains a different digest, execution stops and requires a new run ID.
Stopping stage and report-only inputs are excluded from that digest, so later
targets can be added without falsely defining a new scientific analysis.

Cross-run reuse is powerful only when identity is proven.  A manifest makes
the intended artifact set explicit; hashes prevent path-stable content changes;
and the semantic guard prevents two parameterizations from being written into
one scientific namespace.

\section{Activity quantification}

\subsection{Assay-specific units}

The activity phase requires at least one accepted ATAC and one accepted
H3K27ac library for every context, and at least two accepted libraries of each
assay across the atlas for TMM estimation.

\begin{description}[style=nextline,leftmargin=2.2em]
  \item[Paired-end ATAC]
  Retain proper pairs with \(0 < |\mathrm{TLEN}| < 150\), apply Tn5 offsets,
  and emit both mates as sorted one-base insertion units.

  \item[Paired-end H3K27ac]
  Retain one positive-TLEN record per proper pair and emit the inferred genomic
  fragment once.

  \item[Single-end H3K27ac]
  Convert filtered alignments to BED and extend in the strand direction using
  the manually reviewed fragment length, clipping to chromosome bounds.
\end{description}

Each unit BED has a separately checked total count.  This same unit definition
feeds master-interval counts, 10-kb background counts, H3K27ac summit windows,
and normalized browser coverage.

ATAC accessibility is naturally measured by insertion events, whereas
H3K27ac enrichment is naturally measured by immunoprecipitated fragments.
Using one internally consistent unit per assay prevents mate double-counting
and keeps tables and browser tracks comparable within an assay.

\subsection{Exact variable-width master counts}

For each library, sorted assay units are intersected with the strict,
non-overlapping master registry.  Every master DHS receives a raw count and an
unnormalized CPM/kb value:

\[
  \mathrm{CPM/kb}_{i\ell} =
  \frac{10^9 C_{i\ell}}{N_{\ell} W_i},
\]

where \(C_{i\ell}\) is the raw count for element \(i\) in library \(\ell\),
\(N_{\ell}\) is the total number of usable assay units, and \(W_i\) is the exact
master width in bases.  Coordinates, summit, library, assay, context, raw
count, total units, and CPM/kb are retained together.

\subsection{Assay-specific background TMM}

The reference autosomes are tiled into fixed 10-kb bins, with the terminal bin
on a chromosome clipped to its length.  Raw assay units are counted in every
bin.  For ATAC and H3K27ac separately, bins with zero total count across that
assay are removed and \software{edgeR::calcNormFactors} estimates TMM factors
using the original total filtered assay-unit counts as library sizes.  The
implemented parameters are log-ratio trim 0.30, sum trim 0.05, weighting on,
and \texttt{Acutoff=-1e10}.

For library \(\ell\), the effective library size is

\[
  E_{\ell}=N_{\ell}F_{\ell},
\]

where \(F_{\ell}\) is the assay-specific TMM normalization factor.  The
normalized element signal is

\[
  S_{i\ell}=\frac{10^9 C_{i\ell}}{E_{\ell}W_i}.
\]

The factor table records the assay, context, total units, nonzero background
feature count, factor, effective size, and method.  A software receipt records
the R and edgeR versions and exact parameters.

Estimating compositional factors on master DHS counts risks using the
biological signal of interest as the normalization reference.  Fixed
autosomal background bins provide a broader sampling frame.  Factors remain
assay-specific because insertion counts and ChIP fragments have different
generative processes and should never share a scale factor.

\subsection{Context aggregation and combined activity}

Within each context and assay, normalized biological-library values are
combined by an equal-weight arithmetic mean.  Replicate standard deviations,
raw sums and means, total-unit sums, library IDs, and library counts remain in
the output.  This gives every accepted biological library equal weight
regardless of its number of technical runs.

For context (c) and master DHS (i), the comparative activity score is

\[
  A_{ic}=\sqrt{\overline{S}^{\,\mathrm{ATAC}}_{ic}
                    \;\overline{S}^{\,\mathrm{H3K27ac}}_{ic}}.
\]

This geometric mean is zero when either assay has zero signal and balances two
positive signals without allowing the numerically larger assay to dominate as
strongly as an arithmetic sum.

\section{H3K27ac summit-window signal}

Master intervals are variable width, but H3K27ac may be displaced to either
side of an accessible summit.  The catalog therefore defines three
non-overlapping nominal 500-bp windows around every representative summit:

\begin{align*}
  \text{left}   &: [s-750,\ s-250),\\
  \text{center} &: [s-250,\ s+250),\\
  \text{right}  &: [s+250,\ s+750).
\end{align*}

Windows are clipped at chromosome bounds and their actual widths are retained.
H3K27ac fragments are counted in each window, normalized with the same
library-specific background-TMM effective sizes, and averaged across
biological libraries within the context.  The pipeline then selects the
maximum of the three context means.  It also retains all window-level values
and the mean over the complete 1.5-kb neighborhood.

A single window centered on the Tn5 summit can miss flanking nucleosomes, while
one wide window can dilute a localized signal or merge asymmetric patterns.
The maximum of three fixed local summaries accommodates modest displacement
without changing the master coordinate or optimizing a different window for
every library.

\section{Guarded H3K27ac mixture model}

For each context, model fitting uses only master DHSs that are members of that
context and have positive max-window H3K27ac signal.  Values are transformed
to \(\log_{10}\) scale.  A deterministic expectation--maximization fit compares
one Gaussian with a two-Gaussian mixture.  The lower-mean and higher-mean
components are labeled \texttt{low} and \texttt{high}.

A two-component interpretation is marked supported only when all of the
following guards pass:

\begin{enumerate}
  \item at least 200 positive member DHSs;
  \item \(\mathrm{BIC}_{1}-\mathrm{BIC}_{2} \geq 10\);
  \item Ashman's \(D\geq2\);
  \item each component weight is at least 0.10;
  \item exactly one posterior-probability 0.5 crossing lies between the two
  component means; and
  \item the fitted population has nonzero variance.
\end{enumerate}

If a two-component fit exists, positive member DHSs are assigned using a high
component posterior threshold of 0.5 even when one of the support guards
fails.  Those assignments are retained with
\texttt{mixture\_supported=0},
\texttt{mixture\_guardrail\_warning=1}, and the exact semicolon-separated
failure reasons.  Nonmembers are \texttt{inactive\_in\_context}; accessible
members with zero H3K27ac are
\texttt{accessible\_no\_h3k27ac\_signal}.

The catalog additionally reports a within-context log-signal z score and
descriptive z tiers.  These do not replace the mixture assignment.

A two-Gaussian model will return two components even for many unimodal or
weakly structured distributions.  The guardrails turn model diagnostics into
explicit evidence requirements.  Retaining unsupported assignments with
warnings is preferable to silently discarding data or falsely presenting the
fit as validated.

\section{Regulatory classes and active elements}

The representative summit is compared with the nearest reference TSS on the
same contig.  Ties retain all tied TSS identifiers.  Classes are descriptive:

\begin{center}
\begin{tabular}{lr}
\toprule
Class & Absolute summit--TSS distance \\
\midrule
\texttt{promoter\_associated} & \(\leq250\) bp \\
\texttt{proximal\_enhancer\_like} & 251--1,000 bp \\
\texttt{distal\_enhancer\_like} & \(>1{,}000\) bp \\
\texttt{unclassified\_no\_tss\_on\_contig} & no annotated TSS on contig \\
\bottomrule
\end{tabular}
\end{center}

An active element is a master DHS that is a member of the context and is
assigned to the high H3K27ac mixture component.  The active definition does
not require \texttt{mixture\_supported=1}; unsupported fits remain present but
are visibly flagged in tables, BED names, plots, summaries, and the final
report.  Continuous ATAC, H3K27ac, posterior probability, and combined
activity values are retained for analyses that should not use a binary call.

TSS distance provides a reproducible structural annotation, not proof of
enhancer function or target-gene identity.  Likewise, a high H3K27ac component
is an encyclopedia-style activity label.  The continuous measurements are the
appropriate input for later quantitative scoring and uncertainty-aware work.

\section{Catalog tables, tracks, and IGV sessions}

\subsection{Tabular catalog}

The long catalog contains one master-DHS/context row.  It includes master
coordinates, membership, nearest TSS, regulatory class, normalized ATAC,
all H3K27ac window summaries, mixture diagnostics and posterior, activity
state, and the max-window geometric activity.  The wide catalog contains one
row per master DHS with context-prefixed membership, signal, mixture, warning,
and activity fields.  Each context also receives a compressed table of high
component active elements and a regulatory-class/mixture summary.

\subsection{Portable BED and BigWig tracks}

The catalog creates:

\begin{itemize}
  \item a portable copy of the global master BED;
  \item a context DHS BED containing master coordinates for context-matrix
  members, rather than the wider pooled source-peak boundaries;
  \item an active-element BED9 whose thick one-base interval marks the summit,
  score is the high-component posterior scaled to 0--1,000, item RGB encodes
  regulatory class, and name records a mixture warning; and
  \item one mean ATAC and one mean H3K27ac BigWig for every context.
\end{itemize}

For each accepted library, basewise coverage is scaled by
\(10^6/E_{\ell}\), using the same assay-specific background-TMM effective size
as the tables.  Context tracks are arithmetic means of these scaled library
tracks.  ATAC BigWigs show Tn5 insertion coverage and H3K27ac BigWigs show
fragment coverage.  JSON sidecars record libraries, factors, scale values,
inputs, and output checksums.

\subsection{IGV sessions}

Each context session at
\filep{activity/catalog/igv/<context>.xml} loads five relative resources:
the global master registry first, followed by mean ATAC, mean H3K27ac, context
DHSs, and active elements.  The aggregate
\filep{activity/catalog/all-contexts.igv.xml} loads the master once and the four
context-specific tracks for every context.  Relative paths make the bundle
portable as long as the XML, \filep{bed/}, and \filep{tracks/} remain together.

The standalone \filep{src/build_igv_session.py} builds a compact cross-run view.
The global master DHS registry is first.  Each context then contributes, in
order, its background-TMM mean ATAC Tn5-insertion signal, pooled ATAC qpois
signal, background-TMM mean H3K27ac fragment coverage, context DHSs, and
active cCREs.  The unscaled, depth-dependent pooled MACS3 pileup remains an
auditable master-calling result but is intentionally absent from this compact
comparative session.
Replicate-level H3K27ac signal, H3K27ac peaks, and intermediate ATAC peak
evidence are deliberately omitted.

The tabular catalog is the quantitative record; the browser bundle is a
diagnostic view of the same accepted libraries, units, factors, and coordinates.
Generating both from shared inputs minimizes the risk that a visually
attractive track represents a different normalization or sample selection.

\section{Integrated reporting and provenance}

The report endpoint emits deterministic HTML, PDF, and JSON.  Repeated
\texttt{--report-source-root} arguments freeze a bounded set of upstream
manifests, resolved configurations, QC metrics, TSS plots, fragment
histograms, cross-correlation files, MultiQC reports, and master metrics into
the report configuration.  These sources are checked again by SHA-256 during
report construction.

The integrated report summarizes input/output inventories, BAM and FRiP QC,
ATAC TSS and fragment evidence, ChIP cross-correlation, master statistics, TMM
factors, regulatory classes, mixture fits, exact guardrail warnings, and the
catalog browser assets.  Its JSON sidecar records every declared source and
current output hash, including BEDs, mean BigWigs, and all IGV sessions.

\subsection{Reproducibility mechanisms}

\begin{itemize}
  \item Rule software is isolated in small pinned Conda environments.
  \item Downloads, reused BAMs, master bundles, report sources, and generated
  exports are checksummed.
  \item Resolved configurations preserve derived defaults and provenance.
  \item Temporary outputs are validated and atomically promoted where
  practical.
  \item Deterministic tables use stable ordering; generated gzip files set a
  fixed modification time; unchanged manifests and provenance files are not
  replaced merely to refresh timestamps.
  \item Snakemake targets express completeness and profiles retain
  \texttt{rerun-incomplete: true}.
  \item A changed scientific semantic digest requires a new \texttt{run\_id}.
\end{itemize}

The run-specific \filep{work/<project>/<run_id>/} tree contains reproducible
intermediates rather than archive products.  Large intermediates---trimmed
FASTQs after alignment, lane and marked BAMs, shifted or short-fragment BAMs,
peak-calling insertion BEDs and background-count tables---are declared
temporary and Snakemake removes them after their last consumer.
Quantification assay-unit BEDs are retained below
\filep{activity/quantification/units/}; they are the strict restart contract
for catalog counting and mean-track construction.  The work directory must remain available
while downstream products are being constructed.  After successful completion
it may be removed to clear small files or debris left by interrupted older
runs.  Per-library ChIP CPM BigWigs are not computed because no downstream
result consumes them; pooled ATAC pileup/qpois and context-mean catalog
BigWigs are distinct retained products.  Final replicate-level ATAC peak
evidence is archived below \filep{results/<project>/<run_id>/atac/replicates/}
rather than retained under \filep{work/}.

\section{Principal output inventory}

\begin{longtable}{p{0.24\textwidth} p{0.68\textwidth}}
\toprule
Group & Principal paths below \filep{results/<project>/<run_id>/} \\
\midrule
\endhead
Alignment &
\filep{bam/<library>.final.bam[.bai]};
\filep{qc/alignment/};
\filep{provenance/manifests/final-bams.tsv} \\
QC and review &
\filep{qc/metrics.tsv}, \filep{qc/metrics.json},
\filep{qc/library-review.tsv}, \filep{qc/multiqc/multiqc_report.html},
assay-specific files below \filep{qc/tss/}, \filep{qc/fragments/}, and
\filep{qc/chip/} \\
Context ATAC &
\filep{atac/conditions/<context>/peaks/} and
\filep{atac/conditions/<context>/tracks/} \\
Master bundle &
\filep{atac/master/master_dhs.bed},
\filep{master_dhs_summits.bed},
\filep{master_dhs_membership.tsv},
\filep{master_dhs_context_matrix.tsv}, and
\filep{master_dhs.json} \\
Quantification &
\filep{activity/quantification/units/*.bed.gz},
\filep{activity/quantification/libraries/*.signal.tsv.gz},
\filep{normalization_factors.tsv}, \filep{context_signal.tsv.gz},
\filep{master_dhs_activity.tsv.gz}, metrics and provenance JSON \\
Catalog &
\filep{activity/catalog/master_elements_long.tsv.gz},
\filep{master_elements_wide.tsv.gz}, \filep{active/},
\filep{mixture_models.tsv}, \filep{regulatory_element_summary.tsv}, metrics,
provenance, and mixture plots \\
Browser bundle &
\filep{activity/catalog/bed/}, \filep{activity/catalog/tracks/},
\filep{activity/catalog/igv/<context>.xml}, and
\filep{activity/catalog/all-contexts.igv.xml} \\
Integrated report &
\filep{activity/report/integrated_qc_report.html},
\filep{integrated_qc_report.pdf}, and
\filep{integrated_qc_report.json} \\
Restart contracts &
\filep{provenance/manifests/<stage>.checkpoint.json} at every logical endpoint \\
\bottomrule
\end{longtable}

\section{Representative commands}

\subsection{Accessions through QC}

\begin{lstlisting}
python src/run_pipeline.py resources/atlas_samples_ip_only.tsv \
  --project drosophila-atlas --run-id qc-v1 --genome dm6 \
  --until-stage qc --cores 16
\end{lstlisting}

Copy \filep{qc/library-review.tsv} outside the result namespace, fill the final
review columns, and apply the decisions:

\begin{lstlisting}
python src/review_final_bam_manifest.py \
  results/drosophila-atlas/qc-v1/provenance/manifests/final-bams.tsv \
  --review-table data/reviewed/atlas.library-review.tsv \
  --output data/reviewed/atlas.reviewed.tsv
\end{lstlisting}

\subsection{Reviewed ATAC BAMs through master}

\begin{lstlisting}
python src/run_pipeline.py resources/atlas_samples_ip_only.tsv \
  --from-stage qc \
  --checkpoint-manifest \
    results/drosophila-atlas/qc-v1/provenance/manifests/qc.checkpoint.json \
  --final-bam-manifest data/reviewed/atlas-atac.reviewed.tsv \
  --project drosophila-atlas --run-id master-v1 --genome dm6 \
  --until-stage master --cores 16
\end{lstlisting}

\subsection{Accepted BAMs and master through report}

\begin{lstlisting}
python src/run_pipeline.py resources/atlas_samples_ip_only.tsv \
  --from-stage master \
  --master-manifest \
    results/drosophila-atlas/master-v1/provenance/manifests/master-dhs.tsv \
  --activity-bam-manifest data/reviewed/atac.accepted.tsv \
  --activity-bam-manifest data/reviewed/h3k27ac.accepted.tsv \
  --report-source-root results/drosophila-atlas/master-v1 \
  --report-source-root results/drosophila-h3k27ac/qc-v1 \
  --project drosophila-atlas --run-id catalog-v1 --genome dm6 \
  --until-stage report --cores 16
\end{lstlisting}

Re-running the same command and run ID resumes valid work.  Use
\stage{quantification} or \stage{catalog} as an earlier endpoint.  To advance
later, pass its exported checkpoint as the completed start boundary:

\begin{lstlisting}
python src/run_pipeline.py resources/atlas_samples_ip_only.tsv \
  --from-stage quantification \
  --checkpoint-manifest \
    results/drosophila-atlas/catalog-v1/provenance/manifests/quantification.checkpoint.json \
  --until-stage report --cores 16
\end{lstlisting}

\section{Default parameters that affect interpretation}

\begin{longtable}{p{0.34\textwidth} p{0.18\textwidth} p{0.39\textwidth}}
\toprule
Parameter & Default & Interpretation \\
\midrule
\endhead
Trim quality / minimum length & 20 / 20 bp & Read-level cleanup before alignment \\
Bowtie2 preset / paired maximum & very-sensitive / 2,000 bp & Alignment sensitivity and allowed concordant span \\
MAPQ / duplicate / mitochondrial filtering & 30 / remove / remove & Final biological-library BAM contract \\
ATAC short-fragment maximum & 150 bp, strict & Requires \(0<|\mathrm{TLEN}|<150\) \\
ATAC MACS3 q-value & 0.10 & Lenient candidate stage before geometric refinement \\
qpois exponent range & 2--325 & Progressive signal thresholds \\
qpois retained width & 50--400 bp & Final refined component geometry \\
Context replicate support & 2 libraries, 0.5 coverage & Minimum reproducibility of pooled peaks \\
Master summit span & 150 bp & Complete-linkage limit across context summits \\
Minimum master summit separation & 50 bp & Close disjoint-context clusters may merge below this distance \\
TMM background bins & 10 kb autosomal & Composition-factor feature space \\
H3K27ac windows & three 500-bp windows & Left, center, and right of the master summit \\
Mixture minimum population & 200 positive members & Small contexts cannot support the bimodality claim \\
Mixture evidence & \(\Delta\mathrm{BIC}\geq10\), \(D\geq2\), weights \(\geq0.10\) & Required together with unique crossing and nonzero variance \\
Promoter/proximal boundaries & 250 / 1,000 bp & Absolute summit-to-nearest-TSS distance \\
\bottomrule
\end{longtable}

\section{Pinned principal software}

The rule environments pin the principal tools shown below.  The resolved
Conda environment and report receipts remain the authoritative record for an
individual run.

\begin{center}
\begin{tabular}{ll@{\qquad}ll}
\toprule
Tool & Version & Tool & Version \\
\midrule
Snakemake & 9.23.1 & Bowtie2 & 2.5.5 \\
samtools & 1.23.1 & bedtools & 2.31.1 \\
FastQC & 0.12.1 & Cutadapt & 5.2 \\
MACS3 & 3.0.4 & deepTools & 3.5.6 \\
MultiQC & 1.35 & phantompeakqualtools & 1.2.2 \\
R & 4.5.3 & edgeR & 4.8.2 \\
pyBigWig & 0.3.25 & WeasyPrint & 66.0 \\
\bottomrule
\end{tabular}
\end{center}

\section{Interpretive cautions}

\begin{itemize}
  \item A refined qpois component inherits its candidate call; refinement is
  not a new independent false-discovery-rate procedure.
  \item Replicate support is interval coverage of a pooled peak, not a test of
  identical peak boundaries or identical signal magnitude.
  \item A master DHS represents compatible accessibility evidence across
  contexts.  It is not a fixed-width enhancer, a gene assignment, or a claim
  that the element is active in every context.
  \item Background TMM assumes that a useful reference composition can be
  estimated from autosomal bins.  Factor ranges, library QC, and signal
  distributions should be inspected for strong global or technical shifts.
  \item A high mixture assignment from an unsupported fit is retained for
  transparency but must be interpreted with its warning.  The report exposes
  the exact failed guards.
  \item TSS-proximity classes are annotation categories.  Distal or proximal
  labels do not establish enhancer function, and promoter proximity does not
  prove regulation of the nearest gene.
  \item Browser tracks are normalized context means intended for inspection;
  replicate-level tables and standard deviations remain necessary for
  uncertainty assessment.
  \item ABC scoring and Micro-C integration require downstream workflows that
  consume the continuous signals and frozen master coordinates.
\end{itemize}

\section{Verification of this description}

This document was written against the current sample-sheet schema, launcher,
configuration validators, \filep{workflow/Snakefile}, all included rule files,
the master/consensus/activity/catalog implementations, the README, and the
workflow environment files.  To rebuild the PDF from the repository root:

\begin{lstlisting}
latexmk -pdf -interaction=nonstopmode -halt-on-error \
  -output-directory=docs docs/pipeline-description.tex
\end{lstlisting}

For a scientific run, documentation review is not sufficient verification.
Validate the sample sheet, inspect the resolved configuration and semantic
digest, run an appropriate Snakemake dry-run, review library QC decisions, and
retain the generated manifests and report JSON with the results.

\end{document}
